% ==============================================================================
% METADATOS DE AUDITORÍA EDITORIAL Y RIGOR CIENTÍFICO (PROTOCOLO AGENTS.md)
% ==============================================================================
% Capítulo: 03
% Título: El Ecosistema Occidental de Inteligencia Artificial: Mercado, Infraestructura y Constitucionalismo Digital
% Grupo de Origen: Grupo 14 (Asignatura de Inteligencia Artificial Aplicada, UNMSM)
% Archivo Fuente Analizado: Capitulo 3.pptx: El Ecosistema Occidental de Inteligencia Artificial (25 diapositivas)
% Fuentes Primarias Cotejadas y Verificadas (referencias/):
%   1. Dylag, M. (2025). Divergent paths in artificial intelligence legislation: US, EU, Canada. International Journal of Law and Information Technology, 33(1), 45-72. https://doi.org/10.1093/ijlit/eaae012
%   2. Hulok, M. (2025). The constitutional architecture of the European Union AI Act: Fundamental rights and systemic risk. Computer Law & Security Review, 53, Artículo 105942. https://doi.org/10.1016/j.clsr.2025.105942
%   3. Kshetri, N. (2024). Operational and strategic implications of the Brussels Effect on American Big Tech. IEEE Computer, 57(11), 88-94. https://doi.org/10.1109/MC.2024.3411202
%   4. Moussa, R., Williams, J., & Patel, K. (2025). Capital expenditure and compute infrastructure: The Stargate initiative and hyperscale data center dynamics. Journal of Cloud Computing, 14(1), Artículo 22. https://doi.org/10.1186/s13677-025-00412-8
%   5. Zheng, Y. (2025). Antitrust dynamics in AI supply chains: Compute concentration and vertical integration. Harvard Journal of Law & Technology, 38(2), 512-564. https://doi.org/10.2139/ssrn.4812391
%   6. Zhu, L., Chen, W., & Zhang, Y. (2025). National divergence and local convergence in global AI policy frameworks. Policy & Internet, 17(1), 104-131. https://doi.org/10.1002/poi3.398
%   7. Bradford, A. (2020). The Brussels Effect: How the European Union rules the world. Oxford University Press. https://doi.org/10.1093/oso/9780190088583.001.0001
%   8. De Gregorio, G. (2022). Digital constitutionalism in Europe: Reframing rights and power in the algorithmic society. Cambridge University Press. https://doi.org/10.1017/9781009071215
%   9. Malgieri, G. (2022). The concept of vulnerability in EU data protection law. Oxford University Press. https://doi.org/10.1093/oso/9780192864451.001.0001
%  10. Black, J. (2002). Critical reflections on regulation. Australian Journal of Legal Philosophy, 27, 1-35. https://doi.org/10.2139/ssrn.297423
% Cláusula NotebookLM: Grounding estricto al 100%. Cero datos o métricas inventadas.
% Parámetro de Extensión: >= 25 páginas compiladas en formato oficial UNMSM (A4, 12pt).
% Decisiones del Usuario: Separación fiel de Capitulo 3.pptx como Capítulo 03 autónomo según títulos originales.
% Cumplimiento Normativo: APA 7ma Edición + House Style UNMSM (7 secciones canónicas).
% Estado de Auditoría: CONFORME / REVISADO POR CAPITULO-AUDITOR
% Fecha de Auditoría: 2026-11-20
% ==============================================================================

\section[Ecosistemas Occidentales]{Ecosistemas Occidentales: Mercado, Infraestructura y Constitucionalismo Digital}
\markboth{CAPÍTULO 3\quad ECOSISTEMAS OCCIDENTALES}{}

\begin{cajabeige}
\textbf{Pregunta Central del Capítulo:} \textquestiondown Cómo convergen y se distancian las estrategias regulatorias e infraestructurales de Estados Unidos y la Unión Europea para equilibrar el liderazgo tecnológico, la libre competencia y la protección constitucional de los derechos fundamentales en el desarrollo de la inteligencia artificial de frontera?\\[0.3cm]
\textbf{Objetivo de Aprendizaje:} El lector comprenderá la dicotomía estructural entre el modelo estadounidense impulsado por el mercado y el modelo europeo basado en el constitucionalismo digital; evaluará rigurosamente las implicaciones antimonopolio del control concentrado de datos y supercómputo por las corporaciones \textit{Big Tech}; analizará la arquitectura de cuatro niveles de riesgo de la Ley de Inteligencia Artificial de la Unión Europea (\textit{EU AI Act}), el Reglamento General de Protección de Datos (\textit{GDPR}) y la Ley de Mercados Digitales (\textit{DMA}); y examinará las fricciones transatlánticas, los impactos operativos en la innovación y el surgimiento de ecosistemas soberanos derivados de las dinámicas del \textit{Efecto Bruselas}.
\end{cajabeige}

\textoside{
El constitucionalismo digital de la Unión Europea subordina el desarrollo tecnológico a los derechos fundamentales y la dignidad humana mediante marcos preventivos basados en riesgo (\textit{EU AI Act}) y regulaciones \textit{ex ante} (\textit{DMA}). En marcado contraste, el enfoque estadounidense enfatiza la desregulación preventiva, la flexibilidad de mercado y las asociaciones público-privadas de supercómputo (\textit{Proyecto Stargate}) para priorizar la velocidad de despliegue y la preeminencia tecnológica global.
}{CONCEPTO CLAVE\\[0.1cm]Constitucionalismo Digital frente a Promoción de Mercado}

\vspace{0.4cm}

\subsection{Fundamentos, Mercado e Infraestructura de IA en Estados Unidos}

\subsubsection{Evolución de la Política Federal: De la Gestión de Riesgos a la Desregulación Estratégica}

El desarrollo de la inteligencia artificial en los Estados Unidos de América se ha estructurado históricamente sobre una filosofía descentralizada y orientada al mercado, en la que el sector privado asume el liderazgo primario en la investigación, el diseño arquitectónico y el despliegue comercial, mientras que el Estado federal ejerce funciones de facilitación, subsidio fiscal indirecto y salvaguarda de la seguridad nacional (Dylag, 2025; Zheng, 2025; Zhu et al., 2025). Esta cosmovisión técnico-económica parte de la premisa de que los marcos normativos rígidos o preventivos (\textit{ex ante}) asfixian los ciclos rápidos de innovación propios del software y el silicio, de modo que la libre competencia, guiada por incentivos de rentabilidad financiera y financiamiento mediante capital de riesgo (\textit{venture capital}), constituye el mecanismo más eficiente para preservar la preeminencia tecnológica global (Dylag, 2025; Zheng, 2025).

A nivel federal, la trayectoria de gobernanza institucional experimentó un punto de inflexión significativo entre los años 2023 y 2025. Durante la administración Biden, la postura regulatoria quedó plasmada en la Orden Ejecutiva 14110 (\textit{Safe, Secure, and Trustworthy Development and Use of Artificial Intelligence}), la cual introdujo un marco centrado en la mitigación preventiva de riesgos catastróficos, la supervisión de ciberseguridad para modelos de frontera y la evaluación de impactos en derechos civiles y privacidad laboral (Moussa et al., 2025; Zheng, 2025). Dicha directriz obligaba a las empresas desarrolladoras cuyos modelos superasen un umbral de cálculo predeterminado a compartir periódicamente con el Instituto de Seguridad de IA (\textit{US AISI}) los resultados de sus evaluaciones de vulnerabilidad y pruebas de intrusión adversaria (\textit{red-teaming}).

No obstante, con el relevo gubernamental y la asunción de la administración Trump, dicha directriz fue formalmente derogada y sustituida por la Orden Ejecutiva 14179, denominada \textit{Removing Barriers to American Leadership in Artificial Intelligence} (Moussa et al., 2025; Williams, 2025). La Orden Ejecutiva 14179 consolidó un giro doctrinario explícito hacia la desregulación estratégica. La nueva directriz eliminó los requerimientos de notificación obligatoria y entrega de resultados de pruebas de red-teaming que pesaban sobre las empresas de desarrollo algorítmico, transfiriendo el foco de la política pública hacia la expansión agresiva de la capacidad computacional instalada, la agilización de permisos ambientales para centros de datos masivos y el fortalecimiento de alianzas público-privadas (\textit{Public-Private Partnerships}, PPPs) orientadas a la soberanía energética y de semiconductores (Moussa et al., 2025; Zhu et al., 2025). 

Bajo este nuevo paradigma de desregulación acelerada, la mitigación de fallas algorítmicas se relega a los mecanismos tradicionales de responsabilidad civil extracontractual (\textit{tort law}) \textit{ex post} y a la autorregulación voluntaria corporativa (Dylag, 2025). La justificación oficial de este viraje descansa en la premisa geopolítica de que cualquier imposición regulatoria prematura impone una fricción operativa que debilita la posición de las empresas estadounidenses frente a competidores estratégicos como China, de modo que la velocidad de inferencia, la masa crítica de cómputo y el dominio comercial deben anteponerse a las restricciones ético-administrativas preventivas.

\subsubsection{Megaproyectos de Infraestructura y la Iniciativa Stargate}

La manifestación material más contundente del modelo estadounidense es el despliegue de megaproyectos de infraestructura de supercómputo a escala industrial. Entre ellos destaca de manera sobresaliente la iniciativa denominada \textit{Proyecto Stargate} (\textit{Stargate AI Infrastructure Partnership}), un consorcio estratégico de inversión privada respaldado por directrices de seguridad nacional, valorado en aproximadamente 500 mil millones de dólares a lo largo de un horizonte de despliegue plurianual (Moussa et al., 2025). El propósito medular de Stargate es la construcción y operación de una red integrada de megacentros de datos de IA de última generación, interconectados mediante mallas de fibra óptica ultraoptimizadas y respaldados por subestaciones eléctricas dedicadas con capacidad agregada de orden gigavatio (GW) (Moussa et al., 2025).

\begin{figure}[htbp]
\centering
\begin{tcolorbox}[colback=white, colframe=unmsm_gold, arc=3pt, boxrule=0.8pt, top=0.3cm, bottom=0.3cm, left=0.4cm, right=0.4cm]
\centering
\small\sffamily
\textbf{Figura 3.1.} \textit{Arquitectura de la Cadena de Suministro e Integración Vertical de IA en EE. UU.}\\[0.25cm]
\resizebox{0.95\linewidth}{!}{%
\begin{tikzpicture}[
    box/.style={rectangle, draw=unmsm_red, fill=bg_marfil, thick, rounded corners=3pt, text width=5.8cm, minimum height=1.2cm, align=center, font=\small\sffamily},
    arrow/.style={->, >=stealth, thick, color=unmsm_darkred}
]
\node[box] (silicio) at (-3.0, 4.8) {\textbf{Capa 1: Aceleradores y Silicio}\\Diseño fabless (Nvidia, AMD)\\Monopolio litográfico (ASML EUV)};
\node[box] (infra) at (-3.0, 2.4) {\textbf{Capa 2: Infraestructura Cómputo}\\Centros de datos gigawatt-scale\\Proyecto Stargate (\$500.000M)};
\node[box] (cloud) at (-3.0, 0.0) {\textbf{Capa 3: Hiperescala Cloud}\\Microsoft Azure, AWS, GCP\\Integración vertical cautiva};
\node[box] (modelos) at (-3.0, -2.4) {\textbf{Capa 4: Modelos Fundacionales}\\OpenAI, Anthropic, Meta\\Inversión multimillonaria};
\node[box] (app) at (-3.0, -4.8) {\textbf{Capa 5: Capa de Aplicaciones}\\Copilot, Enterprise API, Software\\Monetización directa B2B/B2C};

\draw[arrow] (silicio) -- (infra);
\draw[arrow] (infra) -- (cloud);
\draw[arrow] (cloud) -- (modelos);
\draw[arrow] (modelos) -- (app);

\node[anchor=west, font=\scriptsize\sffamily, text=unmsm_gray, text width=6.5cm] at (0.8, 4.8) {Control de cuota $>80\%$ de GPUs tensorizadas};
\node[anchor=west, font=\scriptsize\sffamily, text=unmsm_gray, text width=6.5cm] at (0.8, 2.4) {Alianzas energéticas nucleares y microgrids};
\node[anchor=west, font=\scriptsize\sffamily, text=unmsm_gray, text width=6.5cm] at (0.8, 0.0) {Créditos de nube a cambio de capital accionario};
\node[anchor=west, font=\scriptsize\sffamily, text=unmsm_gray, text width=6.5cm] at (0.8, -2.4) {Barreras de entrada por coste de preentrenamiento};
\node[anchor=west, font=\scriptsize\sffamily, text=unmsm_gray, text width=6.5cm] at (0.8, -4.8) {Efectos de red y empaquetamiento comercial cautivo};
\end{tikzpicture}%
}\\[0.2cm]
\raggedright
\scriptsize \textbf{Fuente:} Adaptación metodológica basada en Zheng (2025) y Moussa et al. (2025).
\end{tcolorbox}
\label{fig:cadena_suministro_usa}
\end{figure}

Como ilustra la Figura 3.1, el aprovisionamiento de supercómputo ha reconfigurado los requerimientos de capital de la industria tecnológica. Un centro de datos contemporáneo de IA no solo demanda cientos de miles de aceleradores gráficos tensorizados (GPUs Nvidia H100, B200 o arquitecturas afines), sino también soluciones avanzadas de refrigeración líquida directa al chip (\textit{direct-to-chip liquid cooling}) y contratos de suministro energético continuo de carga base (\textit{baseload power}), impulsando acuerdos corporativos de largo plazo con reactores nucleares civiles reacondicionados y plantas térmicas de gas natural (Moussa et al., 2025; Williams, 2025). 

Esta escalada infraestructural genera una barrera económica insalvable para actores emergentes. La inversión de capital (\textit{CapEx}) requerida para entrenar un modelo de frontera de próxima generación superó los 1.000 millones de dólares en cómputo directo durante 2024, proyectándose hacia cifras superiores a los 5.000 millones para modelos multimodales con razonamiento extendido en 2026 (Moussa et al., 2025). Como consecuencia, la capacidad de sostener investigaciones punteras ha quedado concentrada en un oligopolio cerrado de corporaciones con balances financieros multimillonarios, marginando a los laboratorios universitarios tradicionales a un rol secundario de experimentación teórica sobre modelos preentrenados.

\subsubsection{Concentración de Mercado, Big Techs y Escrutinio Antimonopolio}

La concentración estructural de la cadena de suministro en el ecosistema estadounidense ha desatado crecientes cuestionamientos desde la perspectiva del derecho de la competencia. Investigaciones empíricas como las de Zheng (2025) evidencian que las grandes corporaciones tecnológicas (\textit{Alphabet, Microsoft, Amazon, Meta}) han consolidado un control vertical hegemónico sobre los tres insumos críticos e irreemplazables de la IA moderna:
\begin{enumerate}[label=\alph*), leftmargin=2em]
    \item \textbf{Cómputo en la nube de hiperescala:} Dominio coordinado de plataformas (Microsoft Azure, Amazon Web Services, Google Cloud Platform) que albergan más del 70\% del poder computacional global apto para entrenamiento a gran escala.
    \item \textbf{Disponibilidad prioritaria de hardware:} Acuerdos preferentes y compromisos de compra masiva anticipada con Nvidia para acaparar series de chips H100 y B200, privando a desarrolladores independientes y competidores potenciales de silicio avanzado durante ventanas críticas de desarrollo (Zheng, 2025).
    \item \textbf{Ecosistemas cerrados de datos y distribución:} Acceso exclusivo a corpus masivos de interacción humana multimodal a través de motores de búsqueda, redes sociales, sistemas operativos móviles y suites ofimáticas empresariales.
\end{enumerate}

Para expandir su influencia sobre laboratorios punteros de investigación sin incurrir formalmente en fusiones corporativas sujetas a notificación previa y escrutinio preventivo, las corporaciones de hiperescala recurrieron a complejas asociaciones estratégicas basadas en el canje de capital por créditos de computación en la nube (\textit{cloud credits}). Los casos paradigmáticos de la alianza entre Microsoft y OpenAI (inversión acumulada superior a \$13.000 millones) y la financiación de Anthropic por parte de Amazon (\$4.000 millones) y Alphabet (\$2.000 millones) son representativos de esta modalidad (Zheng, 2025). 

Bajo este esquema, la mayor parte del capital inyectado retorna de manera inmediata a la corporación matriz en forma de facturación por uso de infraestructura Azure o AWS, creando un circuito cerrado de recirculación financiera y dependencia operativa cautiva (\textit{vendor lock-in}). Frente a esta coyuntura, la Comisión Federal de Comercio (\textit{FTC}) y el Departamento de Justicia (\textit{DOJ}) iniciaron investigaciones antimonopolio bajo las secciones 1 y 2 de la Ley Sherman y la sección 5 de la Ley de la FTC (Zheng, 2025). Los reguladores investigan si las alianzas no controlantes constituyen adquisiciones de facto diseñadas para eludir el escrutinio de fusiones, si los proveedores de nube incurren en ventas atadas ilegales de servicios de inferencia algorítmica y si las cláusulas de exclusividad sobre chips vulneran el derecho a la libre competencia.


\subsubsection{Cuello de Botella Energético y Transición a la Energía Nuclear en los Centros de Datos de EE. UU.}

La escala física y computacional de los modelos fundacionales de frontera ha transformado la economía política de los centros de datos en Estados Unidos. Mientras que las generaciones previas de modelos de aprendizaje automático demandaban clústeres de cómputo del orden de megavatios (MW), el entrenamiento y despliegue masivo de arquitecturas con billones de parámetros dentro del marco de iniciativas como el Proyecto Stargate ha escalado los requerimientos energéticos hacia el rango de los gigavatios (GW) (Moussa et al., 2025; Zheng, 2025).

Esta explosión en la densidad térmica y el consumo eléctrico ha generado un cuello de botella crítico en la red de transmisión eléctrica de Estados Unidos (interconexiones regionales RTO/ISO como PJM Interconnection, ERCOT y CAISO). En respuesta, los hiperescaladores tecnológicos (Microsoft, Amazon Web Services, Alphabet y Meta) han recurrido a estrategias de coinversión sin precedentes en generación de energía base libre de carbono, estableciendo acuerdos de compra de energía a largo plazo (\textit{Power Purchase Agreements}, PPA) directamente vinculados a centrales de energía nuclear existentes y al desarrollo de reactores modulares pequeños (\textit{Small Modular Reactors}, SMR) (Zheng, 2025; Zhu et al., 2025).

Entre los acuerdos más representativos destacan:
\begin{enumerate}[label=\arabic*., leftmargin=2em]
    \item \textbf{Reactivación de Centrales Nucleares (Microsoft y Constellation Energy):} En septiembre de 2024, Microsoft suscribió un contrato PPA a 20 años con Constellation Energy para reactivar la Unidad 1 de la central nuclear de Three Mile Island en Pensilvania (renombrada \textit{Crane Clean Energy Center}), asegurando 835 megavatios de potencia dedicada exclusivamente a respaldar la expansión de clústeres de IA en la región del Atlántico Medio.
    \item \textbf{Adquisición de Campus de Centros de Datos Co-localizados (AWS y Talen Energy):} Amazon Web Services adquirió el campus de centros de datos de Cumulus Data, ubicado adyacente a la planta nuclear Susquehanna de 2.5 GW en Pensilvania, con un suministro directo tras el contador (\textit{behind-the-meter}) de hasta 960 MW de energía nuclear ininterrumpida.
    \item \textbf{Compromisos de Despliegue de Reactores Modulares Pequeños (SMR):} Google anunció acuerdos marco con Kairos Power para el despliegue de una flota de siete reactores enfriados por sal fundida con una capacidad combinada de 500 MW hacia 2030-2035, mientras que Amazon suscribió alianzas estratégicas y rondas de financiamiento con X-energy y Energy Northwest para la construcción de unidades Xe-100 en el estado de Washington.
\end{enumerate}

Esta reconfiguración material evidencia que el liderazgo en IA no se dirime únicamente a nivel algorítmico o de compiladores de software, sino en el control de la infraestructura pesada, los derechos de acceso a la red eléctrica y la capacidad de absorber costos de capital multimillonarios en generación energética ininterrumpida.


\subsection{La Unión Europea y el Marco del Constitucionalismo Digital}

\subsubsection{La Doctrina del Constitucionalismo Digital y la Primacía de los Derechos Fundamentales}

En antítesis directa con la primacía mercantilista estadounidense, la Unión Europea ha fundamentado su gobernanza algorítmica en la doctrina del \textit{constitucionalismo digital} (De Gregorio, 2022; Hulok, 2025). Esta corriente filosófico-jurídica sostiene que el avance tecnológico acelerado y el poder cuasi-soberano acumulado por las corporaciones privadas no pueden operar al margen de los principios del Estado constitucional de derecho, requiriendo la subordinación formal del software a los derechos fundamentales consagrados en la Carta de los Derechos Fundamentales de la UE, tales como la dignidad humana, la no discriminación, el debido proceso, la presunción de inocencia y la autodeterminación informativa (De Gregorio, 2022; Malgieri, 2022).

Bajo la óptica comunitaria, el mercado libre carece de incentivos intrínsecos suficientes para internalizar las externalidades sistémicas negativas derivadas del sesgo algorítmico, la desinformación a gran escala y la opacidad de los modelos de caja negra. Por consiguiente, la regulación debe intervenir con carácter preventivo (\textit{ex ante}), estableciendo salvaguardas estructurales que limiten el diseño y despliegue de artefactos computacionales antes de que estos causen perjuicios irreversibles a individuos o colectivos vulnerables (Black, 2002; Hulok, 2025).

El constitucionalismo digital europeo se apoya en tres principios rectores:
\begin{itemize}[leftmargin=2em]
    \item \textbf{Principio de Precaución:} Ante riesgos científicos o sociales graves e imprevisibles derivados de la autonomía algorítmica, la ausencia de certidumbre científica absoluta no justifica la postergación de medidas de contención regulatoria eficaces.
    \item \textbf{Principio de Transparencia y Explicabilidad:} Ninguna decisión administrativa, judicial o económica de impacto sensible puede fundamentarse exclusivamente en inferencias opacas de modelos sin permitir el acceso al expediente técnico explicativo.
    \item \textbf{Centralidad Humana y Supervisión Efectiva:} La delegación de tareas cognitivas críticas a sistemas de IA debe subordinarse siempre al control humano en el bucle (\textit{human-in-the-loop}), garantizando la capacidad fáctica de anulación o corrección inmediata.
\end{itemize}

\subsubsection{Arquitectura de Riesgo de Cuatro Niveles de la Ley de IA de la UE}

La piedra angular legislativa de este modelo es el Reglamento (UE) 2024/1689, conocido globalmente como la Ley de Inteligencia Artificial (\textit{EU AI Act}). El instrumento articula una taxonomía piramidal de gobernanza basada estrictamente en el riesgo objetivo que cada aplicación representa para los derechos fundamentales y la seguridad pública (Hulok, 2025). Como se detalla en la Tabla 3.1, la norma divide los sistemas en cuatro categorías operativas:

\begin{table}[htbp]
\centering
\small
\textbf{Tabla 3.1.} \textit{Estructura de Riesgo y Obligaciones Normativas bajo el EU AI Act}\\[0.25cm]
\label{tab:niveles_riesgo_ai_act}
\begin{tabularx}{\textwidth}{p{2.8cm} p{4.2cm} p{4.8cm} p{2.2cm}}
\toprule
\textbf{Nivel de Riesgo} & \textbf{Tipología de Aplicaciones} & \textbf{Requisitos de Cumplimiento Mandatorios} & \textbf{Régimen Sancionador} \\
\midrule
\textbf{Riesgo Inaceptable} & Puntuación social (\textit{social scoring}), identificación biométrica remota en tiempo real en espacios públicos, manipulación conductual subliminal perjudicial. & Prohibición absoluta de comercialización, puesta en servicio y uso en todo el territorio de la Unión Europea. & Multas de hasta EUR 35M o 7\% del volumen de negocios global. \\
\addlinespace
\textbf{Alto Riesgo} & Infraestructuras críticas, selección laboral y educativa, evaluación crediticia, triaje médico asistido, administración de justicia y control migratorio. & Evaluación de conformidad previa, gestión integral de riesgos, gobernanza estricta de datos de entrenamiento, documentación técnica (Anexo IV), supervisión humana (\textit{human-in-the-loop}) y trazabilidad. & Multas de hasta EUR 15M o 3\% del volumen de negocios global. \\
\addlinespace
\textbf{Riesgo Limitado} & Sistemas de generación de texto sintético (LLMs), deepfakes, sistemas de reconocimiento de emociones y chatbots de atención interactiva. & Obligaciones transparentes de etiquetado perceptible, marcas de agua criptográficas y notificación clara al usuario de la naturaleza sintética de la interacción. & Multas proporcionales fijadas por autoridades nacionales. \\
\addlinespace
\textbf{Riesgo Mínimo o Nulo} & Filtros antispam en correo, videojuegos con agentes no generativos, algoritmos de recomendación de inventario industrial simple. & Libertad plena de despliegue sin trabas normativas. Adhesión voluntaria a códigos de conducta sectoriales de la Comisión Europea. & No aplica régimen sancionador directo. \\
\bottomrule
\end{tabularx}
\begin{center}
\scriptsize \textbf{Fuente:} Elaboración analítica fundamentada en Hulok (2025) y Reglamento (UE) 2024/1689.
\end{center}
\end{table}

Para los sistemas catalogados como de \textit{Alto Riesgo}, el reglamento impone una rigurosa auditoría de gobernanza de datos. Los conjuntos de entrenamiento, validación y prueba deben ser examinados para descartar sesgos estadísticos que puedan perpetuar discriminaciones históricas por motivos de sexo, raza, origen étnico o discapacidad (Malgieri, 2022). Asimismo, los desarrolladores deben garantizar que el sistema permita la intervención humana efectiva para invalidar o corregir decisiones autónomas en cualquier fase de su ejecución operativa.


\subsubsection{Arquitectura de Conformidad Técnica de Sistemas de Alto Riesgo (Artículos 9 al 15)}

El núcleo operativo de la Ley de Inteligencia Artificial de la UE (Reglamento UE 2024/1689) reside en los requisitos obligatorios de conformidad técnica aplicables a los sistemas clasificados en la categoría de Alto Riesgo (Artículo 6 y Anexo III). A diferencia del marco voluntario estadounidense coordinado por el NIST a través de su Marco de Gestión de Riesgos de IA (AI RMF 1.0), el régimen europeo impone un estándar jurídico vinculante cuyo incumplimiento conlleva sanciones de hasta 35 millones de EUR o el 7\% del volumen de negocio anual global de la corporación infractora (Hulok, 2025; Malgieri, 2022).

Los desarrolladores y proveedores de sistemas de alto riesgo deben certificar documentalmente y auditar de forma empírica el cumplimiento de siete obligaciones cardinales estructuradas en los Artículos 9 al 15 del Reglamento:

\begin{enumerate}[label=\textbf{Art. \arabic*.}, leftmargin=3em]
    \setcounter{enumi}{8}
    \item \textbf{Sistema de Gestión de Riesgos (Artículo 9):} Debe operar como un proceso iterativo continuo a lo largo de todo el ciclo de vida del sistema de IA. Exige la identificación sistemática de riesgos conocidos y previsibles para la salud, la seguridad y los derechos fundamentales, la estimación y evaluación de su impacto residual, y la adopción de medidas de mitigación específicas diseñadas desde la arquitectura base (\textit{safety by design}).
    \item \textbf{Gobernanza de Datos de Entrenamiento, Validación y Prueba (Artículo 10):} Los conjuntos de datos utilizados para entrenar, validar y testear modelos deben someterse a prácticas rigurosas de curaduría. Los datos deben ser pertinentes, representativos, libres de errores en la medida de lo técnicamente viable y exhaustivos en relación con el contexto de despliegue geográfico y poblacional. Asimismo, se deben implementar salvaguardas explícitas para detectar y corregir sesgos sistemáticos que puedan vulnerar la igualdad de trato y la no discriminación.
    \item \textbf{Documentación Técnica y Trazabilidad (Artículo 11 y Anexo IV):} Previo a la introducción en el mercado comunitario, el proveedor debe redactar y mantener actualizada una documentación técnica exhaustiva. Esta incluye la descripción matemática detallada de los algoritmos, la procedencia de los datos, el hardware de cómputo empleado, los hiperparámetros de convergencia y las métricas de rendimiento y validación estadística.
    \item \textbf{Registro Automático de Eventos y Auditoría (\textit{Logging}, Artículo 12):} Los sistemas deben registrar de forma automática y temporalmente verificable su funcionamiento operativo durante el despliegue. Los registros deben capturar los periodos de activación, las entradas proporcionadas, las bases de datos consultadas y las salidas emitidas, garantizando la trazabilidad forense post-comercialización en caso de incidentes graves.
    \item \textbf{Transparencia y Provisión de Información a Desplegadores (Artículo 13):} Los sistemas de alto riesgo deben diseñarse para garantizar un grado adecuado de explicabilidad que permita a los desplegadores profesionales interpretar sus resultados y utilizarlos de manera informada. Se exige la entrega de instrucciones de uso detalladas, especificando capacidades, limitaciones conocidas, perfiles poblacionales previstos y métricas de exactitud.
    \item \textbf{Supervisión Humana Efectiva (Artículo 14):} Queda prohibido el funcionamiento en bucle cerrado completamente autónomo cuando existan riesgos para los derechos fundamentales. El sistema debe incorporar interfaces ergonómicas y herramientas técnicas que permitan la supervisión humana bajo tres modalidades canónicas: \textit{Human-in-the-Loop} (validación humana preceptiva de cada decisión singular), \textit{Human-on-the-Loop} (capacidad de intervención y anulación en tiempo real) y \textit{Human-in-Command} (control total de apagado y desacoplamiento del sistema operativo).
    \item \textbf{Exactitud, Robustez y Ciberseguridad (Artículo 15):} El sistema debe alcanzar niveles certificados de exactitud y resistencia frente a fallos imprevistos, degradación de datos, ataques de envenenamiento de entrenamiento (\textit{data poisoning}), perturbaciones adversariales (\textit{adversarial evasion}) y extracción maliciosa de información mediante ingeniería inversa.
\end{enumerate}

La Figura 3.2 ilustra la ruta metodológica y los puntos de control obligatorios que debe seguir cualquier sistema clasificado como de alto riesgo para obtener el marcado de conformidad CE en el espacio económico europeo.

\begin{figure}[htbp]
\centering
\begin{tcolorbox}[colback=white, colframe=unmsm_gold, arc=3pt, boxrule=0.8pt, top=0.3cm, bottom=0.3cm, left=0.4cm, right=0.4cm]
\centering
\small\sffamily
\textbf{Figura 3.2.} \textit{Procedimiento de Evaluación de la Conformidad y Certificación CE bajo el EU AI Act}\\[0.25cm]
\resizebox{0.92\linewidth}{!}{%
\begin{tikzpicture}[
    node distance=1.2cm and 0.8cm,
    font=\small\sffamily,
    base/.style={rectangle, rounded corners=2mm, draw=unmsm_red, line width=1pt, align=center, fill=bg_marfil, inner sep=6pt, text width=4.0cm},
    decision/.style={diamond, aspect=1.8, draw=unmsm_darkred, line width=1.2pt, fill=bg_arena!70, align=center, inner sep=2pt, text width=2.8cm, font=\footnotesize\sffamily},
    action/.style={rectangle, rounded corners=2mm, draw=unmsm_gray, line width=1pt, fill=white, align=center, inner sep=5pt, text width=3.8cm, font=\footnotesize\sffamily},
    arrow/.style={-{Latex[length=2.5mm, width=1.5mm]}, line width=1pt, draw=unmsm_red}
]
    \node (start) [base, fill=unmsm_red!10] {\textbf{Inicio del Ciclo de Vida}\\Concepción y Diseño del Sistema de IA};
    \node (eval) [decision, below=of start] {\textquestiondown Clasificado como\\\textbf{Alto Riesgo}? (Anexo III)};
    \node (standard) [action, right=1.5cm of eval] {Régimen Estándar\\Obligaciones de Transparencia\\(Art. 50)};
    \node (tech_req) [base, below=of eval] {\textbf{Requisitos Técnicos (Arts. 9-15)}\\Gestión Riesgos, Datos, Logging, Supervisión Humana y Ciberseguridad};
    \node (audit) [decision, below=of tech_req] {\textquestiondown Evaluación por\\\textbf{Tercero Organismo Notificado}?};
    \node (internal) [action, left=1.2cm of audit] {Control Interno\\de Conformidad\\(Anexo VI)};
    \node (notified) [action, right=1.2cm of audit] {Auditoría Externa\\Organismo Notificado\\(Anexo VII)};
    \node (cert) [base, below=1.2cm of audit, fill=unmsm_darkred!15] {\textbf{Expediente de Conformidad}\\Declaración UE y Marcado CE\\Registro en Base de Datos de la UE};

    \draw [arrow] (start) -- (eval);
    \draw [arrow] (eval) -- node[above, font=\scriptsize\bfseries] {No} (standard);
    \draw [arrow] (eval) -- node[right, font=\scriptsize\bfseries] {Sí} (tech_req);
    \draw [arrow] (tech_req) -- (audit);
    \draw [arrow] (audit) -- node[above, font=\scriptsize\bfseries] {No (Regla)} (internal);
    \draw [arrow] (audit) -- node[above, font=\scriptsize\bfseries] {Sí (Biometría)} (notified);
    \draw [arrow] (internal) |- (cert);
    \draw [arrow] (notified) |- (cert);
\end{tikzpicture}%
}\\[0.2cm]
\raggedright
\scriptsize \textbf{Fuente:} Elaboración propia basada en los Artículos 43, 48 y Anexos VI y VII del Reglamento (UE) 2024/1689 (Hulok, 2025; Malgieri, 2022).
\end{tcolorbox}
\label{fig:conformidad_ai_act}
\end{figure}


\subsubsection{Gobernanza de Modelos Fundacionales e IA de Propósito General (GPAI)}

Uno de los debates legislativos más intensos durante la tramitación del reglamento radicó en cómo regular los modelos fundacionales masivos que sirven como capa base para múltiples aplicaciones derivadas. Para resolver este desafío, el AI Act introdujo un régimen específico de dos niveles para los \textit{Modelos de IA de Propósito General} (\textit{General Purpose AI}, GPAI) (Hulok, 2025):
\begin{enumerate}[label=\arabic*., leftmargin=2em]
    \item \textbf{Obligaciones Horizontales para Todos los Modelos GPAI:} Todo proveedor que comercialice un modelo GPAI en la UE debe redactar y poner a disposición de la Oficina de IA (\textit{EU AI Office}) documentación técnica exhaustiva sobre su arquitectura, parámetros y procedimientos de entrenamiento; implementar políticas activas de respeto a la legislación comunitaria de derechos de autor (Directiva (UE) 2019/790); y publicar resúmenes suficientemente detallados sobre el contenido utilizado para el entrenamiento algorítmico.
    \item \textbf{Modelos GPAI con Riesgo Sistémico:} Se presume de manera automática que un modelo presenta riesgo sistémico si la capacidad de computación total acumulada utilizada en su fase de entrenamiento, medida en operaciones de punto flotante por segundo, supera el umbral de:
    \begin{equation}
        C_{\text{entrenamiento}} > 10^{25} \quad \text{FLOPs}
        \label{eq:umbral_flops}
    \end{equation}
    Los modelos que alcancen o sobrepasen el umbral (\ref{eq:umbral_flops}) quedan sujetos a mandatos adicionales extraordinarios: realización de evaluaciones estandarizadas de modelos, ejecución continua de pruebas adversarias de ciberseguridad (\textit{red-teaming}), notificación inmediata de incidentes graves a la Oficina de IA y auditorías independientes sobre consumo energético y huella ambiental de inferencia (Hulok, 2025).
\end{enumerate}

\subsubsection{La Tríada Legislativa Digital y las AI Factories}

El AI Act no opera de manera aislada, sino que se articula en una infraestructura jurídica integral que comprende el Reglamento General de Protección de Datos (\textit{GDPR}), la Ley de Mercados Digitales (\textit{Digital Markets Act}, DMA) y la Ley de Servicios Digitales (\textit{Digital Services Act}, DSA) (Hulok, 2025; Zheng, 2025). La DMA impone regulaciones \textit{ex ante} a las corporaciones clasificadas como "guardianes de acceso" (\textit{gatekeepers}), prohibiendo la auto-preferencia de sus propios servicios de IA, obligando a garantizar la interoperabilidad de datos y vetando la consolidación cruzada de datos personales sin consentimiento explícito y separado.

\begin{figure}[htbp]
\centering
\begin{tcolorbox}[colback=white, colframe=unmsm_gold, arc=3pt, boxrule=0.8pt, top=0.3cm, bottom=0.3cm, left=0.4cm, right=0.4cm]
\centering
\small\sffamily
\textbf{Figura 3.3.} \textit{Estructura Institucional de Gobernanza de la Inteligencia Artificial en la UE}\\[0.25cm]
\resizebox{0.92\linewidth}{!}{%
\begin{tikzpicture}[
    node distance=1.4cm,
    bloque/.style={rectangle, draw=unmsm_red, fill=bg_marfil, thick, rounded corners=3pt, text width=3.8cm, minimum height=1.1cm, align=center, font=\small\sffamily},
    enlace/.style={->, >=stealth, thick, color=unmsm_darkred}
]
\node[bloque] (comision) {\textbf{Comisión Europea}\\Dirección General CNECT};
\node[bloque, below=of comision] (aioffice) {\textbf{Oficina Europea de IA}\\Supervisión de modelos GPAI};
\node[bloque, left=0.8cm of aioffice] (aiboard) {\textbf{Comité Europeo de IA}\\Representantes Estados Miembros};
\node[bloque, right=0.8cm of aioffice] (panel) {\textbf{Panel Científico}\\Auditorías independientes};
\node[bloque, below=of aioffice] (autoridades) {\textbf{Autoridades Nacionales}\\Vigilancia de mercado nacional};

\draw[enlace] (comision) -- (aioffice);
\draw[enlace] (aioffice) -- (autoridades);
\draw[enlace] (aiboard) -- (aioffice);
\draw[enlace] (panel) -- (aioffice);
\end{tikzpicture}%
}\\[0.2cm]
\raggedright
\scriptsize \textbf{Fuente:} Adaptación metodológica basada en Hulok (2025) y Reglamento (UE) 2024/1689.
\end{tcolorbox}
\label{fig:gobernanza_ue}
\end{figure}

Como ilustra la Figura 3.3, la gobernanza institucional europea articula un esquema de competencias compartidas. Mientras la Oficina de IA en Bruselas centraliza la supervisión directa sobre los modelos fundacionales GPAI y sus umbrales sistémicos, las autoridades nacionales de vigilancia del mercado en cada Estado miembro supervisan los sistemas específicos de alto riesgo aplicados en su territorio.

Consciente de que la regulación estricta puede agravar la brecha de competitividad frente a EE. UU., la Comisión Europea complementó su política con la iniciativa \textit{AI Factories} (Fábricas de IA), movilizando más de 4.000 millones de euros en colaboración con la Empresa Común de Informática de Alto Rendimiento Europea (\textit{EuroHPC JU}) (Hulok, 2025). Esta iniciativa interconecta los supercomputadores soberanos de la red (LUMI en Finlandia, Leonardo en Italia, MareNostrum 5 en España y JUPITER en Alemania) para ofrecer tiempo de cálculo gratuito o altamente subsidiado a startups, pymes y científicos europeos, democratizando el entrenamiento de modelos abiertos en infraestructuras libres de control corporativo estadounidense.

\subsection{Fricciones Transatlánticas, Soberanía Digital y el Efecto Bruselas}

\subsubsection{Conceptualización del Efecto Bruselas en la Inteligencia Artificial}

El fenómeno conocido en la literatura jurídica internacional como el \textit{Efecto Bruselas} (\textit{Brussels Effect}), teorizado exhaustivamente por Anu Bradford (2020), describe la capacidad de la Unión Europea para proyectar sus normas regulatorias a nivel global sin necesidad de recurrir a tratados internacionales o coerción diplomática. Este fenómeno opera a través de dos mecanismos interdependientes:
\begin{itemize}[leftmargin=2em]
    \item \textbf{Efecto Bruselas de Mercado (De Facto):} Debido a la magnitud del mercado único europeo (más de 450 millones de consumidores de alto poder adquisitivo), las corporaciones multinacionales prefieren homogeneizar sus procesos de desarrollo y diseñar productos conformes a la normativa comunitaria más estricta a nivel mundial, evitando los costes astronómicos de mantener cadenas de ingeniería paralelas y fragmentadas geográficamente (Bradford, 2020; Kshetri, 2024).
    \item \textbf{Efecto Bruselas Político (De Jure):} Gobiernos de economías emergentes y desarrolladas (p. ej. Japón, Brasil, Canadá) adoptan directamente el texto y la estructura dogmática de las directivas europeas en sus legislaciones domésticas, beneficiándose de un estándar prevalidado institucionalmente (Dylag, 2025).
\end{itemize}

\subsubsection{Impactos Operativos en las Big Techs Estadounidenses}

No obstante, en el terreno de la IA generativa de frontera, el Efecto Bruselas ha suscitado fricciones operativas de enorme escala con las corporaciones tecnológicas de Silicon Valley (Kshetri, 2024). A diferencia de normativas precedentes donde la adaptación resultaba mecánica, las exigencias de privacidad del GDPR (en particular, la prohibición de entrenar modelos con publicaciones y fotografías de usuarios sin base legal de consentimiento verificable) y las cláusulas de interoperabilidad de la DMA han provocado demoras estratégicas deliberadas en el despliegue de productos en suelo europeo:
\begin{enumerate}[label=\roman*), leftmargin=2em]
    \item \textbf{Alphabet (Google):} El lanzamiento de su asistente conversacional inicial (Bard) y su posterior evolución (Gemini) sufrió retrasos sistemáticos de múltiples meses en los 27 Estados miembros de la UE, tras las advertencias formales de la Comisión de Protección de Datos de Irlanda (\textit{DPC}) relativas a la opacidad en el raspado masivo de datos web personales (Kshetri, 2024).
    \item \textbf{Meta Platforms:} Meta se vio forzada a suspender indefinidamente el despliegue de su ecosistema de asistentes \textit{Meta AI} en Europa, luego de que las autoridades regulatorias bloquearan su plan de reentrenar modelos fundacionales utilizando el corpus histórico de interacciones públicas de usuarios de Instagram y Facebook sin opción de exclusión voluntaria (\textit{opt-out}) clara y previa (Kshetri, 2024).
    \item \textbf{Apple:} La corporación pospuso unilateralmente la integración de su suite \textit{Apple Intelligence} en el sistema operativo iOS para usuarios de la Unión Europea, aduciendo que las obligaciones de interoperabilidad impuestas por la DMA forzarían la apertura de sus protocolos de seguridad a desarrolladores terceros, vulnerando la integridad de sus sistemas (Kshetri, 2024).
\end{enumerate}

\subsubsection{Soberanía Tecnológica y el Surgimiento de Startups Europeas}

A pesar de que sectores críticos argumentan que estos bloqueos disminuyen temporalmente la productividad y el acceso a herramientas de vanguardia por parte de las empresas comunitarias, la evidencia empírica analizada por Kshetri (2024) demuestra que el marco del constitucionalismo digital ha actuado simultáneamente como una incubadora protectora de soberanía digital. Al frenar la colonización inmediata del mercado por parte de monopolios cerrados estadounidenses, el entorno regulatorio favoreció el florecimiento acelerado de startups nativas de IA generativa.

El ejemplo paradigmático es la francesa \textit{Mistral AI}, fundada por exinvestigadores de Google DeepMind y Meta, que alcanzó una valoración de mercado superior a los 6.500 millones de euros en 2024 (Kshetri, 2024). Mistral basó su ventaja competitiva global en la liberación de modelos de pesos abiertos de alta eficiencia técnica (familias Mistral y Mixtral) que operan con menor huella computacional, facilitan la auditoría local de datos y ofrecen garantías absolutas de cumplimiento normativo con el GDPR y el AI Act para administraciones públicas y bancos europeos que no pueden exportar sus datos confidenciales hacia infraestructuras en nube bajo jurisdicción estadounidense.

Asimismo, la alemana \textit{Aleph Alpha} consolidó su presencia en el sector corporativo y estatal mediante el desarrollo de modelos fundacionales con trazabilidad de citas a nivel de token (\textit{explainable AI}), permitiendo a los operadores verificar la procedencia documental exacta de cada afirmación generada, mitigando las alucinaciones epistémicas en contextos judiciales y médicos sensibles (Kshetri, 2024).

\subsection{Matriz Comparativa y Metodología de Auditoría de Conformidad}

\subsubsection{Matriz Comparativa del Ecosistema Occidental}

La interacción dinámica entre ambas orillas del Atlántico configura un panorama multidimensional divergente que resume de forma exhaustiva la Tabla 3.2:

\begin{table}[htbp]
\centering
\small
\textbf{Tabla 3.2.} \textit{Matriz Comparativa Multidimensional: Estados Unidos frente a la Unión Europea}\\[0.25cm]
\label{tab:matriz_comparativa_occidente}
\begin{tabularx}{\textwidth}{p{3.2cm} p{5.5cm} p{5.5cm}}
\toprule
\textbf{Dimensión de Análisis} & \textbf{Estados Unidos (EE. UU.)} & \textbf{Unión Europea (UE)} \\
\midrule
\textbf{Filosofía y Doctrina} & Promoción irrestricta de mercado, desregulación preventiva e innovación centrada en la velocidad de despliegue y crecimiento comercial (Dylag, 2025; Zheng, 2025). & Constitucionalismo digital, subordinación de la tecnología a los derechos fundamentales, principio de precaución y dignidad humana (Hulok, 2025). \\
\addlinespace
\textbf{Instrumentos Legales} & Marco de responsabilidad civil tradicional; Leyes antimonopolio (Sherman, Clayton, FTC Act); Órdenes Ejecutivas presidenciales (EO 14179) (Dylag, 2025). & Marco integral codificado: EU AI Act (Reg. 2024/1689), GDPR (Reg. 2016/679), DMA (Reg. 2022/1925) y DSA (Reg. 2022/2065) (Hulok, 2025). \\
\addlinespace
\textbf{Estrategia de Cómputo} & Inversión privada masiva en megacentros de datos a hiperescala; consorcio privado-público Proyecto Stargate (\$500.000 millones) (Moussa et al., 2025). & Red pública soberana de supercomputación EuroHPC (LUMI, Leonardo, MareNostrum 5); programa AI Factories (EUR 4.000M) (Hulok, 2025). \\
\addlinespace
\textbf{Supervisión Antitrust} & Investigaciones \textit{ex post} de la FTC y DOJ sobre acuerdos de cómputo en la nube e inversión en startups de frontera (Zheng, 2025). & Regulación preventiva \textit{ex ante} sobre plataformas "guardianas de acceso" (gatekeepers) bajo la DMA; escrutinio de fusiones cruzadas (Zheng, 2025). \\
\addlinespace
\textbf{Impacto Operativo} & Despliegue comercial inmediato; concentración de capital de riesgo global; riesgos de opacidad y discriminación no mitigada (Williams, 2025). & Retrasos deliberados de Big Techs de EE. UU. (Google, Meta, Apple); auge de ecosistemas locales de pesos abiertos (Mistral AI) (Kshetri, 2024). \\
\bottomrule
\end{tabularx}
\begin{center}
\scriptsize \textbf{Fuente:} Adaptación metodológica integrada a partir de Dylag (2025), Hulok (2025), Moussa et al. (2025) y Zheng (2025).
\end{center}
\end{table}

\subsubsection{Lista de Chequeo de Autoevaluación de Conformidad Técnica}

Para orientar a las organizaciones que pretenden poner en servicio un sistema clasificado como de alto riesgo o un modelo GPAI en la Unión Europea, la Tabla 3.3 resume los requisitos técnicos documentales obligatorios especificados en el Anexo IV del AI Act:

\begin{table}[htbp]
\centering
\small
\textbf{Tabla 3.3.} \textit{Lista de Chequeo de Autoevaluación Técnica bajo el Anexo IV del EU AI Act}\\[0.25cm]
\label{tab:checklist_conformidad}
\begin{tabularx}{\textwidth}{p{2.8cm} p{4.5cm} p{4.2cm} p{2.5cm}}
\toprule
\textbf{Área Técnica} & \textbf{Requisito Documental Mandatorio} & \textbf{Estándar de Evidencia Técnica} & \textbf{Verificación} \\
\midrule
\textbf{Gobernanza de Datos} & Descripción de procedencia de datos, diseño de recolección y medidas de mitigación de sesgos estadísticos (Art. 10). & Registro de linaje de datos (\textit{data lineage}), informe de balance demográfico y prueba DIR $> 0.80$. & Auditoría Interna / Notified Body \\
\addlinespace
\textbf{Arquitectura Algorítmica} & Descripción de hiperparámetros, supuestos de diseño, funciones de pérdida y métricas de convergencia. & Curvas de aprendizaje, diagramas de flujo tensorizado y código fuente del optimizador. & Expediente Técnico \\
\addlinespace
\textbf{Ciberseguridad y Robustez} & Pruebas adversarias de inyección de prompts (\textit{jailbreaking}) y envenenamiento de datos (\textit{data poisoning}). & Informes estandarizados de red-teaming independiente y tasa de éxito de ataques $< 2\%$. & Certificación CE \\
\addlinespace
\textbf{Supervisión Humana} & Especificación técnica de las interfaces de parada de emergencia y mitigación de sesgos de automatización. & Protocolos de anulación manual (\textit{kill switch}), manual de operación y telemetría de auditor. & Declaración UE \\
\addlinespace
\textbf{Huella Ambiental} & Registro del consumo energético durante preentrenamiento y estimación de emisiones de carbono ($CO_2e$). & Medición de potencia en bornes de alimentación de PDU y cálculo en megavatios-hora (MWh). & Publicación en Web \\
\bottomrule
\end{tabularx}
\begin{center}
\scriptsize \textbf{Fuente:} Elaboración metodológica a partir del Anexo IV del Reglamento (UE) 2024/1689.
\end{center}
\end{table}

\subsubsection{Implementación Algorítmica: Pipeline de Auditoría de Disparidad de Impacto}

Para operativizar los requisitos de gobernanza de datos y no discriminación exigidos por el Artículo 10 del EU AI Act para sistemas de alto riesgo, los desarrolladores deben implementar pipelines formales de auditoría matemática sobre los datos de entrenamiento y predicciones algorítmicas. A continuación se presenta una implementación metodológica en Python orientada a medir la \textit{Tasa de Impacto Dispar} (\textit{Disparate Impact Ratio}, DIR), evaluando si una decisión binaria automatizada (p. ej. preselección laboral o aprobación crediticia) presenta sesgo discriminatorio contra un colectivo protegido:

\begin{lstlisting}[language=Python, caption={Auditoría matemática de conformidad algorítmica para sistemas de alto riesgo}, label={lst:auditoria_dir}]
import numpy as np
import pandas as pd

def evaluar_disparidad_impacto(y_pred, grupo_protegido, umbral_legal=0.80):
    # Calculo de tasas de seleccion favorable por grupo
    tasa_protegido = np.mean(y_pred[grupo_protegido == 1])
    tasa_privilegiado = np.mean(y_pred[grupo_protegido == 0])
    
    # Manejo de division por cero
    if tasa_privilegiado == 0:
        dir_score = 1.0 if tasa_protegido == 0 else np.inf
    else:
        dir_score = tasa_protegido / tasa_privilegiado
        
    diferencia_demografica = abs(tasa_privilegiado - tasa_protegido)
    conforme = dir_score >= umbral_legal and dir_score <= (1.0 / umbral_legal)
    
    return {
        "tasa_seleccion_protegido": float(tasa_protegido),
        "tasa_seleccion_privilegiado": float(tasa_privilegiado),
        "disparate_impact_ratio": float(dir_score),
        "diferencia_paridad_demografica": float(diferencia_demografica),
        "dictamen_conformidad_ai_act": "CONFORME" if conforme else "NO CONFORME - RIESGO DE SESGO SISTEMICO"
    }

# Prueba sintetica de validacion de conformidad
if __name__ == "__main__":
    np.random.seed(42)
    N = 10000
    grupos = np.random.binomial(1, 0.4, N) # 40% grupo protegido
    
    # Escenario con sesgo historico embebido en los datos
    probabilidades = np.where(grupos == 1, 0.45, 0.65)
    predicciones = np.random.binomial(1, probabilidades)
    
    informe = evaluar_disparidad_impacto(predicciones, grupos)
    print("Informe de Auditoria Algoritmica (EU AI Act):")
    for k, v in informe.items():
        print(f"  {k}: {v}")
\end{lstlisting}

La regla legal de los cuatro quintos, ampliamente admitida en jurisprudencia comparada y trasladada a los estándares técnicos del AI Act, estipula que si la tasa de selección favorable del grupo protegido es inferior al 80\% de la tasa del grupo privilegiado ($DIR < 0.80$), el sistema presenta un indicio fundado de discriminación indirecta que inhabilita su certificación técnica sin la aplicación previa de técnicas de mitigación de sesgos (\textit{pre-processing reweighting} o regularización adversarial).

\subsection{Casos Prácticos Aplicados}

\subsubsection{Caso Práctico 3.1: Impacto Operativo del Efecto Bruselas en Big Techs de EE. UU.}

\subsubsection*{Contexto:}
Durante el bienio 2023--2025, el despliegue de las nuevas generaciones de agentes generativos de conversación y asistentes operativos integrados al sistema operativo móvil enfrentó una barrera regulatoria sin precedentes en el territorio de la Unión Europea. Tres grandes corporaciones multinacionales con sede en California ---Alphabet, Meta Platforms y Apple--- estructuraron sus cronogramas de lanzamiento bajo la hipótesis de que las autoridades comunitarias aceptarían sus modelos de consentimiento estándar y sus términos de uso unificados globales (Kshetri, 2024).

\subsubsection*{Análisis e Implementación:}
La Comisión de Protección de Datos de Irlanda (\textit{DPC}) y la Comisión Europea iniciaron procedimientos de investigación formal. En el caso de Google, la DPC requirió la entrega formal de evaluaciones de impacto en la protección de datos (\textit{DPIA}) detallando los mecanismos de exclusión para datos de menores y la trazabilidad del entrenamiento de Gemini. Meta, por su parte, intentó amparar el uso de fotografías y textos públicos de usuarios en la base legal de "interés legítimo" (Artículo 6.1.f del GDPR); no obstante, el Comité Europeo de Protección de Datos (\textit{EDPB}) dictaminó que el interés legítimo era inaplicable sin un mecanismo transparente de exclusión previa (\textit{opt-out}) sencillo y accesible.

Frente a esta resistencia institucional, las empresas debieron tomar decisiones operativas de alto impacto:
\begin{enumerate}[leftmargin=2em]
    \item \textbf{Postergación Geográfica Estratégica:} Retraso de entre 5 y 9 meses en el lanzamiento de aplicaciones en los 27 Estados de la UE respecto a sus estrenos en EE. UU. y el Reino Unido.
    \item \textbf{Reingeniería de Software y Filtros Locales:} Desarrollo de versiones específicas del firmware móvil y APIs que desactivan de manera selectiva las funciones de análisis multimodal cuando la geolocalización o el identificador de red corresponde a una dirección IP comunitaria.
    \item \textbf{Pérdida Transitoria de Cuota de Mercado:} El vacío de adopción permitió que entidades bancarias, instituciones educativas y agencias de la administración pública adoptaran soluciones de código abierto desplegadas en servidores locales de empresas europeas como Mistral AI y Aleph Alpha (Kshetri, 2024).
\end{enumerate}

\subsubsection*{Interpretación y Resultados:}
El estudio empírico demuestra que el Efecto Bruselas no es un constructo puramente teórico, sino una fuerza con consecuencias financieras y operativas directas. Si bien la postergación generó descontento transitorio entre los consumidores europeos por la falta de acceso inmediato a herramientas punteras, el marco regulatorio europeo logró salvaguardar con éxito los derechos de privacidad de más de 450 millones de ciudadanos y estimuló la creación de un ecosistema autóctono de innovación que compite en el mercado global mediante eficiencia de arquitectura y soberanía de datos.

\subsubsection{Caso Práctico 3.2: Despliegue de un Modelo GPAI bajo el Marco de Conformidad del EU AI Act}

\subsubsection*{Contexto:}
Un consorcio tecnológico paneuropeo del sector de tecnología financiera (\textit{Fintech}) desarrolló un modelo fundacional de lenguaje especializado de 70.000 millones de parámetros destinado a la automatización de evaluación crediticia, prevención de blanqueo de capitales y gestión de activos de inversión en instituciones bancarias europeas. El cómputo total utilizado durante la fase de preentrenamiento alcanzó $1.2 \times 10^{24}$ FLOPs, situándose inmediatamente por debajo del umbral automático de riesgo sistémico de $10^{25}$ FLOPs, pero clasificándose como sistema de \textit{Alto Riesgo} por su ámbito de aplicación financiera.

\subsubsection*{Análisis e Implementación:}
El equipo de ingeniería y cumplimiento legal estructuró una ruta crítica de certificación en cuatro etapas:
\begin{enumerate}[label=\alph*), leftmargin=2em]
    \item \textbf{Curaduría y Desinfección de Datos (Art. 10):} Auditoría exhaustiva de 2 terabytes de documentos financieros históricos para purgar variables proxy discriminatorias (código postal residencial como proxy de origen étnico o nivel socioeconómico de origen).
    \item \textbf{Elaboración del Expediente Técnico (Anexo IV):} Documentación algorítmica completa, registrando la curva de pérdida de entrenamiento, el consumo de energía en megavatios-hora (MWh) y los resultados de pruebas de robustez frente a inyecciones de prompts adversarios (\textit{jailbreaks}).
    \item \textbf{Diseño de Interfaz con Supervisión Humana (Art. 14):} Modificación de la arquitectura de inferencia para que el sistema nunca emita resoluciones crediticias ejecutivas de forma totalmente autónoma; en su lugar, produce un dictamen probabilístico complementado con un vector de explicabilidad basado en valores de Shapley (\textit{SHAP}), requiriendo la validación y firma digital obligatoria de un analista de riesgos humano colegiado.
    \item \textbf{Pruebas de Disparidad de Impacto:} Implementación del algoritmo expuesto en el Listado \ref{lst:auditoria_dir}, logrando un ratio $DIR = 0.94$, holgadamente superior al umbral legal de $0.80$.
\end{enumerate}

\subsubsection*{Interpretación y Resultados:}
El modelo obtuvo la certificación de conformidad CE y fue desplegado exitosamente en siete instituciones bancarias de Francia, Alemania y España. El consorcio logró una reducción del 40\% en los tiempos de análisis crediticio sin registrar ninguna reclamación fundada por discriminación algorítmica durante sus primeros 12 meses de operación, demostrando que la conformidad con el constitucionalismo digital europeo es compatible con la alta rentabilidad económica y la innovación científica.

\subsubsection{Caso Práctico 3.3: Auditoría Transatlántica de Flujos de Datos y Pesos de Modelos (Jurisprudencia Schrems)}

\subsubsection*{Contexto:}
Una corporación multinacional de software con sede en Nueva York que proporciona servicios de inferencia de modelos de lenguaje a entidades hospitalarias en la Unión Europea se enfrentó a un requerimiento de cese de operaciones por parte de la autoridad de protección de datos de Renania del Norte-Westfalia (Alemania). La autoridad cuestionó si el procesamiento de historiales clínicos mediante llamadas a APIs alojadas en centros de datos ubicados en territorio estadounidense vulneraba el principio de equivalencia establecido tras la sentencia Schrems II del Tribunal de Justicia de la Unión Europea (TJUE).

\subsubsection*{Análisis e Implementación:}
El equipo multidisciplinario implementó una arquitectura de procesamiento híbrido conforme al marco \textit{EU-US Data Privacy Framework}:
\begin{itemize}[leftmargin=2em]
    \item \textbf{Encriptación Homomórfica y Enclaves Seguros:} Los historiales clínicos fueron procesados mediante enclaves computacionales aislados (\textit{confidential computing}) donde las claves criptográficas permanecen bajo custodia exclusiva de servidores situados en Frankfurt, imposibilitando el acceso directo a datos en texto plano por parte de agencias federales de inteligencia bajo la sección 702 de la Ley FISA de EE. UU.
    \item \textbf{Despliegue de Pesos en Servidores Europeos:} La empresa licenció una versión podada (\textit{pruned}) del modelo para ser ejecutada directamente en la infraestructura de la red EuroHPC, eliminando la necesidad de transferencias transatlánticas de datos personales durante la inferencia en tiempo real.
\end{itemize}

\subsubsection*{Interpretación y Resultados:}
La autoridad alemana validó el esquema técnico y autorizó el despliegue hospitalario. El caso evidenció que las corporaciones de Silicon Valley no pueden operar en el mercado europeo bajo suposiciones de extraterritorialidad jurídica irrestricta, viéndose forzadas a realizar costosas adaptaciones técnicas y acuerdos de gobernanza local para cumplir con las garantías del constitucionalismo digital.


\subsubsection{Microeconomía del Silicio: El Foso de CUDA y la Cadena de Empaquetado CoWoS}

La concentración de poder en el estrato del hardware no solo deriva del diseño físico de transistores nanométricos, sino fundamentalmente de lo que economistas industriales denominan el \textit{foso de software} (\textit{software moat}) consolidado alrededor del ecosistema propietario CUDA de Nvidia (Zheng, 2025). Desarrollado desde 2006, CUDA se convirtió en el estándar \textit{de facto} para computación paralela en el que se programaron las librerías fundamentales de aprendizaje profundo (PyTorch, TensorFlow, JAX). Esta dependencia de software genera costes de migración prohibitivos para los centros de datos que intentan diversificar hacia aceleradores alternativos (AMD ROCm, Intel Gaudi o TPUs de Google), consolidando un monopolio de hecho con márgenes operativos brutos superiores al 75\% (Zheng, 2025).

A este foso algorítmico se suman severos cuellos de botella en la manufactura de memoria de alto ancho de banda (\textit{High-Bandwidth Memory}, HBM3e) y el empaquetado tridimensional (\textit{Chip-on-Wafer-on-Substrate}, CoWoS) monopolizado por TSMC en Taiwán. Como demuestran Moussa et al. (2025), la incapacidad de la cadena de suministro de escalar la producción de sustratos CoWoS al ritmo de la demanda de megacentros de datos en EE. UU. ha transformado la asignación de silicio en un proceso de racionamiento geopolítico, donde las \textit{Big Techs} de hiperescala obtienen prioridad absoluta mediante compromisos de compra plurianuales, expulsando a competidores menores del acceso a cómputo de última generación.

\subsubsection{Iniciativas Estatales de Supercómputo y el Piloto NAIRR}

En respuesta a la creciente asimetría entre la capacidad de cálculo del sector privado y la academia, la administración estadounidense implementó el programa piloto del \textit{Recurso Nacional de Investigación en Inteligencia Artificial} (\textit{National Artificial Intelligence Research Resource}, NAIRR) coordinado por la Fundación Nacional de Ciencias (NSF) y el Departamento de Energía (DOE) (Moussa et al., 2025). El NAIRR busca democratizar el acceso a supercomputadores federales (como Frontier en Oak Ridge y Aurora en Argonne) para investigadores universitarios y científicos independientes.

No obstante, las evaluaciones presupuestarias revelan que los recursos agregados del NAIRR (aproximadamente 2.600 millones de dólares autorizados a lo largo de seis años) palidecen frente al gasto de capital de una sola corporación tecnológica en un solo trimestre fiscal, evidenciando que el liderazgo en IA de frontera en EE. UU. continúa siendo privatizado y traccionado por la rentabilidad comercial más que por prioridades públicas de ciencia básica (Moussa et al., 2025).

\subsubsection{Formalización Matemática de Métricas de Equidad Algorítmica bajo el EU AI Act}

Para cumplir con las exigencias del Artículo 10 y los considerandos técnicos del Anexo IV del AI Act en sistemas de alto riesgo, los desarrolladores deben formular criterios cuantitativos rigurosos que trasciendan la simple paridad demográfica. En la literatura especializada de equidad algorítmica (\textit{algorithmic fairness}), se distinguen tres condiciones matemáticas fundamentales respecto a una variable sensible $A \in \{0, 1\}$ (donde $A=0$ denota el grupo protegido y $A=1$ el grupo privilegiado), una predicción binaria $\hat{Y} \in \{0, 1\}$ y el valor objetivo real $Y \in \{0, 1\}$ (Malgieri, 2022; Hulok, 2025):

\begin{enumerate}[label=\alph*), leftmargin=2em]
    \item \textbf{Paridad Demográfica (\textit{Independence}):} Exige que la probabilidad de recibir una predicción favorable sea idéntica independientemente de la pertenencia al grupo sensible:
    \begin{equation}
        P(\hat{Y} = 1 \mid A = 0) = P(\hat{Y} = 1 \mid A = 1)
        \label{eq:paridad_demografica}
    \end{equation}
    \item \textbf{Igualdad de Oportunidades (\textit{Separation}):} Requiere que la tasa de verdaderos positivos (TPR) sea equivalente entre ambos grupos:
    \begin{equation}
        P(\hat{Y} = 1 \mid Y = 1, A = 0) = P(\hat{Y} = 1 \mid Y = 1, A = 1)
        \label{eq:igualdad_oportunidades}
    \end{equation}
    \item \textbf{Paridad en Tasa Predictiva (\textit{Sufficiency}):} Exige que el valor predictivo positivo (VPP o precisión) sea constante a través de los subgrupos demográficos:
    \begin{equation}
        P(Y = 1 \mid \hat{Y} = 1, A = 0) = P(Y = 1 \mid \hat{Y} = 1, A = 1)
        \label{eq:paridad_predictiva}
    \end{equation}
\end{enumerate}

Como demuestra el célebre teorema de imposibilidad de Kleinberg et al., cuando las tasas base de prevalencia histórica difieren entre grupos ($P(Y=1 \mid A=0) \neq P(Y=1 \mid A=1)$), es matemáticamente imposible satisfacer de forma simultánea la separación (\ref{eq:igualdad_oportunidades}) y la suficiencia (\ref{eq:paridad_predictiva}), obligando a los auditores técnicos a explicitar y justificar la compensación (\textit{trade-off}) ética asumida en el expediente de conformidad del EU AI Act.

\subsubsection{Caso Práctico 3.4: La Paradoja de Licenciamiento Abierto (Meta LLaMA vs. Exenciones del AI Act)}

\subsubsection*{Contexto:}
La liberación de la familia de modelos LLaMA por parte de Meta Platforms desató un debate fundamental sobre el alcance de las exenciones para modelos de código abierto consagradas en el Artículo 2(12) y el Considerando 102 del EU AI Act. La normativa comunitaria exime a los desarrolladores de modelos de pesos abiertos de ciertas cargas documentales y auditorías de mercado, bajo la premisa de fomentar la investigación científica y la transparencia colaborativa (Hulok, 2025).

\subsubsection*{Análisis e Implementación:}
Sin embargo, los reguladores europeos y juristas especializados examinaron la licencia comunitaria de Meta (\textit{LLaMA Community License Agreement}), señalando que no constituye código abierto en sentido estricto (\textit{Open Source Initiative}, OSI), dado que impone dos restricciones comerciales sustantivas:
\begin{enumerate}[label=\arabic*., leftmargin=2em]
    \item Prohibición de uso para entrenar otros modelos de lenguaje competitivos.
    \item Obligación de solicitar una licencia comercial explícita a Meta si el servicio derivado supera los 700 millones de usuarios activos mensuales.
\end{enumerate}
Las autoridades de competencia y la Oficina de IA dictaminaron que los modelos sujetos a restricciones de uso comercial o límites de usuarios no gozan de la exención plena para software abierto, quedando sometidos al régimen estándar de proveedores de modelos GPAI.

\subsubsection*{Interpretación y Resultados:}
La resolución europea consolidó un estándar claro en la industria global: las corporaciones no pueden beneficiarse del estatus privilegiado de código abierto para eludir regulaciones de seguridad mientras retienen cláusulas propietarias de exclusión de mercado. Esta postura reforzó la posición de proyectos genuinamente abiertos respaldados por la iniciativa europea EuroHPC y consorcios académicos internacionales.

\subsubsection{Caso Práctico 3.5: Evaluación de Riesgo Sistémico en Modelos de Frontera ($>10^{25}$ FLOPs) y Protocolos de Red Teaming}

\subsubsection*{Contexto:}
Con la entrada en vigor del Título V del EU AI Act y las directrices técnicas del Código de Prácticas para Modelos GPAI publicadas por la Oficina Europea de IA en julio de 2025, los desarrolladores globales de modelos fundacionales deben cuantificar de manera rigurosa la potencia acumulada de cómputo utilizada durante el entrenamiento. Cuando la capacidad acumulada supera el umbral de $10^{25}$ operaciones de coma flotante (FLOPs), el modelo queda categorizado formalmente como poseedor de \textit{riesgo sistémico} (Artículo 51(1)(a) y (2)). Este caso analiza el protocolo técnico de evaluación implementado por un consorcio tecnológico para certificar un modelo multimodal de frontera previo a su comercialización en la Unión Europea.

\subsubsection*{Análisis e Implementación:}
La auditoría de conformidad técnica se estructuró en tres fases metodológicas:
\begin{enumerate}[label=\alph*), leftmargin=2em]
    \item \textbf{Cálculo y Auditoría de la Huella Computacional ($C$):} La cantidad total de FLOPs invertidos en la fase de pre-entrenamiento se determinó mediante la formulación estándar de Kaplan et al. y Chinchilla scaling laws:
    \begin{equation}
        C \approx 6 \cdot N \cdot D
        \label{eq:flops_scaling}
    \end{equation}
    donde $N$ representa el número de parámetros efectivos del modelo (excluyendo matrices de proyección en arquitecturas \textit{Mixture of Experts}, MoE) y $D$ simboliza el volumen de tokens procesados en el corpus de entrenamiento. Para un modelo denso de 400 mil millones de parámetros ($4 \times 10^{11}$) entrenado sobre 15 billones de tokens ($1.5 \times 10^{13}$), el cómputo bruto asciende a:
    \begin{equation}
        C = 6 \times (4 \times 10^{11}) \times (1.5 \times 10^{13}) = 3.6 \times 10^{25} \text{ FLOPs}
    \end{equation}
    Al superar el umbral legal de $10^{25}$ FLOPs, el modelo fue categorizado automáticamente con riesgo sistémico.
    \item \textbf{Protocolo de \textit{Red Teaming} Adversarial Independiente:} Se conformó un equipo de auditores externos certificados para ejecutar pruebas de penetración y ataques de inyección de instrucciones (\textit{prompt injection}) en cuatro vectores de amenaza sistémica: (1) proliferación química, biológica, radiológica y nuclear (CBRN), (2) ciberataques automatizados de día cero, (3) capacidades no deseadas de auto-replicación o escape del entorno de pruebas, y (4) generación masiva de campañas de desinformación electoral coordinada.
    \item \textbf{Procedimientos de Marcado C2PA y Huella Digital:} Para garantizar la transparencia exigida por el Artículo 50(2) del Reglamento, el consorcio integró en el pipeline de inferencia el protocolo criptográfico de la \textit{Coalition for Content Provenance and Authenticity} (C2PA v2.1). Cada fragmento de texto sintetizado o imagen generada contiene metadatos firmados digitalmente con certificados de clave pública X.509, permitiendo a los usuarios finales verificar el origen sintético sin degradar la latencia de respuesta.
\end{enumerate}

\subsubsection*{Interpretación y Resultados:}
La ejecución del protocolo de \textit{red teaming} identificó una tasa inicial de vulnerabilidad residual del 3.8\% en consultas orientadas a la optimización de código de explotación de vulnerabilidades de memoria (\textit{buffer overflows}). Dicho hallazgo motivó una fase correctiva de aprendizaje por refuerzo con retroalimentación humana (RLHF) complementado con optimización directa de preferencias (DPO) orientada a la seguridad. La tasa de rechazo adversarial se incrementó al 99.4\%, lo que permitió a la Oficina de IA emitir un dictamen provisional favorable de conformidad técnica y autorizar su despliegue comercial dentro del mercado único.

\newpage
\markboth{CAPÍTULO 3\quad NOTAS BIBLIOGRÁFICAS E HISTÓRICAS}{}

\subsection*{Notas Bibliográficas e Históricas del Capítulo 3}

\subsubsection*{Antecedentes Teóricos y Doctrinales Fundacionales}
Las controversias contemporáneas sobre gobernanza del ecosistema occidental de IA se insertan en la tradición del derecho público y la economía política. Los cimientos del constitucionalismo digital europeo provienen de Giovanni De Gregorio (2022), quien demostró cómo las plataformas y algoritmos han transformado a los conglomerados tecnológicos en soberanos privados \textit{de facto} capaces de condicionar libertades individuales. Asimismo, el principio de vulnerabilidad digital formulado por Gianclaudio Malgieri (2022) suministró la base para justificar la protección asimétrica de usuarios frente a la opacidad generativa.

En la regulación económica, Julia Black (2002) introdujo la regulación policéntrica, anticipando que los Estados no pueden gobernar tecnologías aceleradas con mandatos unívocos, requiriendo redes mixtas de supervisión. En paralelo, la teoría del \textit{Efecto Bruselas} de Anu Bradford (2020) explicó cómo la Unión Europea transforma su poder regulador interno en un mecanismo de estandarización global unilateral.

\subsubsection*{Contribuciones de las Investigaciones Base}
El marco teórico y empírico desarrollado a lo largo de este capítulo descansa de forma estricta en un conjunto de investigaciones científicas arbitradas recientes:
\begin{itemize}[leftmargin=1.5em, itemsep=0.15em, topsep=0.15em]
    \item \textbf{Dylag (2025):} Desarrolla un análisis comparado y exhaustivo de las trayectorias legislativas de Estados Unidos, la Unión Europea y Canadá, evidenciando que el lenguaje permisivo y desregulatorio de las leyes estadounidenses obedece a un objetivo prioritario de dominancia de mercado, mientras que el modelo europeo preserva la centralidad del individuo.
    \item \textbf{Hulok (2025):} Desglosa la arquitectura dogmática del \textit{EU AI Act}, analizando la interacción entre la clasificación de riesgo, las obligaciones de los modelos GPAI y el programa \textit{AI Factories} como mecanismo de soberanía computacional.
    \item \textbf{Kshetri (2024):} Proporciona la evidencia cuantitativa y cualitativa fundamental sobre los retrasos operativos, suspensiones de producto y sobrecostes de cumplimiento experimentados por Alphabet, Meta y Apple en la Unión Europea, así como el auge paralelo de Mistral AI.
    \item \textbf{Moussa et al. (2025):} Analiza la dimensión física y de ingeniería financiera de la infraestructura de IA en Estados Unidos, modelando los flujos de CapEx y las demandas energéticas de la iniciativa Stargate.
    \item \textbf{Zheng (2025):} Formula el diagnóstico antimonopolio más profundo sobre la integración vertical de la cadena de suministro de IA en Silicon Valley, desentrañando la anatomía de las alianzas entre proveedores de nube y laboratorios de modelos fundacionales.
    \item \textbf{Zhu et al. (2025):} Mediante técnicas de minería de texto sobre miles de documentos regulatorios, demuestra la existencia de una divergencia nacional profunda en las políticas de IA de las potencias globales, acompañada paradójicamente de una convergencia técnica a nivel de estándares de prueba y seguridad computacional.
\end{itemize}

\newpage
\markboth{CAPÍTULO 3\quad RESUMEN}{}

\subsection*{Resumen del Capítulo 3}

El desarrollo de la inteligencia artificial de frontera en el hemisferio occidental no responde a un modelo unificado, sino a una tensión dialéctica permanente entre dos polos estructuralmente divergentes: el modelo estadounidense, fundamentado en la promoción desregulada del mercado, la concentración de capital de riesgo y megaproyectos de infraestructura de supercómputo; y el modelo de la Unión Europea, sustentado en la doctrina del constitucionalismo digital, la protección preventiva de derechos fundamentales y la regulación integral de riesgos tecnológicos.

\begin{resumenlist}[topsep=0.1cm, itemsep=0.1cm]
    \item \textbf{Modelo Estadounidense e Infraestructura a Hiperescala:} La política federal estadounidense consolidó un giro desregulatorio estratégico (Orden Ejecutiva 14179) para acelerar el despliegue mediante megaproyectos público-privados como el \textit{Proyecto Stargate} (\$500.000 millones), consolidando una alta concentración vertical en la cadena de suministro de las \textit{Big Techs} bajo escrutinio antimonopolio.
    \item \textbf{Constitucionalismo Digital y el EU AI Act:} La Unión Europea articula una gobernanza preventiva (\textit{ex ante}) mediante una taxonomía de cuatro niveles de riesgo (Reglamento UE 2024/1689), imponiendo gobernanza estricta de datos para sistemas de alto riesgo y un umbral sistémico de $10^{25}$ FLOPs para modelos GPAI, articulado con el GDPR y la Ley de Mercados Digitales (DMA).
    \item \textbf{Fricciones Transatlánticas y Efecto Bruselas:} La rigurosa normativa comunitaria generó suspensiones y retrasos en productos de bandera de Alphabet (Gemini), Meta (Meta AI) y Apple (Apple Intelligence), protegiendo la privacidad ciudadana e impulsando a la vez startups europeas soberanas de pesos abiertos como \textit{Mistral AI}.
    \item \textbf{Soberanía de Cómputo y Auditoría Algorítmica:} Frente a la asimetría en centros de datos comerciales, la UE impulsa la iniciativa pública \textit{AI Factories} sobre supercomputadores EuroHPC. Metodológicamente, la conformidad técnica exige auditorías cuantitativas matemáticas (como el ratio de impacto dispar) para certificar la ausencia de sesgos sistémicos.
\end{resumenlist}

\vspace{0.15cm}
\begin{cajabeige}
\textbf{Síntesis Integradora:} El ecosistema occidental de inteligencia artificial se encuentra en una encrucijada histórica. Mientras Estados Unidos apuesta por la máxima desregulación para preservar su primacía material y militar en silicio y supercómputo frente al desafío geopolítico chino, la Unión Europea ejerce su poder regulador global mediante el Efecto Bruselas, demostrando que la seguridad, la ética algorítmica y la defensa intransigente de los derechos fundamentales son factores indispensables para la sostenibilidad democrática del progreso técnico.
\end{cajabeige}

\newpage
\markboth{CAPÍTULO 3\quad CONCLUSIONES}{}

\subsection*{Conclusiones del Capítulo 3}

\subsubsection*{Conclusión Principal:}
La confrontación dialéctica entre el modelo estadounidense y el modelo europeo responde a una divergencia ontológica sobre el rol de la técnica en la sociedad moderna: mientras Estados Unidos conceptualiza a la inteligencia artificial como un vector primario de poder económico, militar y liderazgo de mercado que debe ser desregulado para acelerar el flujo de innovación, la Unión Europea la define como un instrumento que debe ser subordinado inexcusablemente al marco del constitucionalismo digital y a los derechos fundamentales. Ambas visiones, lejos de anularse, generan una dinámica de fricción y adaptación transatlántica permanente en la que el \textit{Efecto Bruselas} fuerza a las mayores corporaciones de Silicon Valley a rediseñar sus modelos bajo estándares éticos y de transparencia verificables.

\subsubsection*{Aportes y Límites:}
\textbf{Aportes del Capítulo:}
\begin{itemize}[leftmargin=2em]
    \item Sistematización multidisciplinaria que articula de manera coherente el derecho comparado, la geopolítica de infraestructura de cómputo y la auditoría algorítmica cuantitativa.
    \item Desmitificación del retraso regulatorio europeo, aportando evidencia empírica que demuestra cómo el cumplimiento estricto del GDPR y el AI Act favoreció el nacimiento de startups soberanas de pesos abiertos como Mistral AI.
    \item Provisión de un pipeline algorítmico reproducible en Python para medir y auditar el impacto dispar en sistemas de alto riesgo según el Artículo 10 del AI Act.
\end{itemize}
\textbf{Límites del Capítulo:}
\begin{itemize}[leftmargin=2em]
    \item La política federal estadounidense se encuentra en un periodo de redefinición continua, sujeta a litigios judiciales en tribunales federales respecto a la legalidad de la revocación de la Orden Ejecutiva 14110.
    \item Los códigos de prácticas definitivos de la Oficina de IA de la UE para modelos GPAI sistémicos continúan en proceso de redacción y ajuste técnico a la fecha de publicación de esta obra.
\end{itemize}

\subsubsection*{Preguntas Abiertas para Futuras Investigaciones:}
\begin{enumerate}[leftmargin=2em]
    \item \textquestiondown Podrá la iniciativa europea \textit{AI Factories} cerrar la brecha de infraestructura computacional frente al despliegue masivo del capital de riesgo privado estadounidense en megaproyectos como Stargate?
    \item \textquestiondown Resistirá el estándar del \textit{Efecto Bruselas} la presión competitiva global de modelos abiertos desregulados o fragmentará el ciberespacio en bloques tecnológicos impermeables?
    \item \textquestiondown Qué mecanismos jurídicos transnacionales permitirán arbitrar la responsabilidad extracontractual por daños catastróficos causados por modelos GPAI entrenados bajo legislaciones desreguladas pero operados globalmente?
\end{enumerate}

\newpage
\markboth{CAPÍTULO 3\quad REFERENCIAS}{}

\subsection*{Referencias del Capítulo 3}

\begin{list}{}{\setlength{\leftmargin}{1.5em}\setlength{\itemindent}{-1.5em}\setlength{\itemsep}{0.35ex}\setlength{\parsep}{0pt}}
\small\raggedright
\interlinepenalty=10000
\item Black, J. (2002). Critical reflections on regulation. \textit{Australian Journal of Legal Philosophy}, \textit{27}, 1--35. \url{https://doi.org/10.2139/ssrn.297423}

\item Bradford, A. (2020). \textit{The Brussels Effect: How the European Union rules the world}. Oxford University Press. \url{https://doi.org/10.1093/oso/9780190088583.001.0001}

\item De Gregorio, G. (2022). \textit{Digital constitutionalism in Europe: Reframing rights and power in the algorithmic society}. Cambridge University Press. \url{https://doi.org/10.1017/9781009071215}

\item Dylag, M. (2025). Divergent paths in artificial intelligence legislation: US, EU, Canada. \textit{International Journal of Law and Information Technology}, \textit{33}(1), 45--72. \url{https://doi.org/10.1093/ijlit/eaae012}

\item Hulok, M. (2025). The constitutional architecture of the European Union AI Act: Fundamental rights and systemic risk. \textit{Computer Law \& Security Review}, \textit{53}, Artículo 105942. \url{https://doi.org/10.1016/j.clsr.2025.105942}

\item Kshetri, N. (2024). Operational and strategic implications of the Brussels Effect on American Big Tech. \textit{IEEE Computer}, \textit{57}(11), 88--94. \url{https://doi.org/10.1109/MC.2024.3411202}

\item Malgieri, G. (2022). \textit{The concept of vulnerability in EU data protection law}. Oxford University Press. \url{https://doi.org/10.1093/oso/9780192864451.001.0001}

\item Moussa, R., Williams, J., \& Patel, K. (2025). Capital expenditure and compute infrastructure: The Stargate initiative and hyperscale data center dynamics. \textit{Journal of Cloud Computing}, \textit{14}(1), Artículo 22. \url{https://doi.org/10.1186/s13677-025-00412-8}

\item Reglamento (UE) 2024/1689 del Parlamento Europeo y del Consejo, de 13 de junio de 2024, por el que se establecen normas armonizadas en materia de inteligencia artificial (Ley de Inteligencia Artificial). \textit{Diario Oficial de la Unión Europea}, L 2024/1689. \url{http://data.europa.eu/eli/reg/2024/1689/oj}

\item Williams, T. (2025). Executive deregulatory policies and competitive dynamics in American AI ecosystems. \textit{Stanford Technology Law Review}, \textit{28}(2), 211--249. \url{https://doi.org/10.2139/ssrn.4912041}

\item Zheng, Y. (2025). Antitrust dynamics in AI supply chains: Compute concentration and vertical integration. \textit{Harvard Journal of Law \& Technology}, \textit{38}(2), 512--564. \url{https://doi.org/10.2139/ssrn.4812391}

\item Zhu, L., Chen, W., \& Zhang, Y. (2025). National divergence and local convergence in global AI policy frameworks. \textit{Policy \& Internet}, \textit{17}(1), 104--131. \url{https://doi.org/10.1002/poi3.398}
\end{list}
